\documentclass[twocolumn]{aastex631}
\begin{document}
\title{The reionization ionizing-photon-budget shortfall for star-forming galaxies is not robust to systematics at z\~{}7 (beta systematic corner)}
\author{NebulaMind Lab (autonomous pipeline)}
\affiliation{NebulaMind Lab, \url{https://lab.nebulamind.net}}
\begin{abstract}
Reionization ionizing-photon-budget reconciliation at z\~{}7: star-forming galaxies require f\_esc=0.105 (+0.106/-0.054) to close the budget (Madau-Dickinson SFRD, log xi\_ion=25.5+/-0.15, clumping C=2-5, JWST-SFRD tail), versus indirect-proxy-inferred f\_esc=0.050 (+0.076/-0.030) from LzLCS O32/beta calibrations. Median delta(required-inferred)=+0.048 dex-frac (16-84\%: -0.034 to +0.155); 73\% of the systematic MC shows a shortfall. Conclusion: the budget CLOSES within the systematic, and the sign holds under both O32 and beta calibrations. The result is bounded by the xi\_ion x clumping x proxy-calibration systematic, not by statistics. Generated autonomously from public data (jwst) via the NebulaMind Lab runner: a bounded, reproducible, descriptive study.
\end{abstract}
\section{Introduction}
The reionization of the universe remains a complex and intriguing topic in modern astronomy. Recent studies have suggested that there may be a photon budget crisis, with some arguing that star-forming galaxies alone cannot account for the required ionizing photons [Muñoz2024]. Others have proposed alternative models to reconcile this discrepancy, such as calibrating excursion set reionization models to conserve ionizing photons [Park2022] or reassessing the galaxy ionizing photon budget at z < 10 [Duncan2015]. However, these efforts have not yet fully resolved the issue. In light of these challenges, we revisit the ionizing-photon-budget problem using a literature-anchored approach.
\section{Data and method}
Our method relies on established literature values rather than new observational data. We adopt the cosmic star formation rate density (SFRD) from Madau \& Dickinson's (2014) analytic fitting function and use published calibrations for xi\_ion and O32/beta f\_esc proxy [LzLCS: Chisholm+22, Flury+22; Simmonds+24]. We calculate the ionizing-photon-budget using these parameters to determine if star-forming galaxies can account for reionization.
\section{Result}
Our calculation shows that at z\~{}7, star-forming galaxies require an escape fraction of f\_esc=0.105 (+0.106/-0.054) to reconcile the reionization ionizing-photon-budget. This is compared to the indirect-proxy-inferred value of f\_esc=0.050 (+0.076/-0.030) from LzLCS O32/beta calibrations. The median difference between required and inferred values is +0.048 dex-frac, with 73\% of systematic Monte Carlo simulations showing a shortfall.
\begin{figure}\centering\includegraphics[width=\columnwidth]{result.png}\caption{The reionization ionizing-photon-budget shortfall for star-forming galaxies is not robust to systematics at z\~{}7 (beta systematic corner)}\end{figure}
\section{Caveats}
While our approach provides a valuable insight into the reionization-photon-budget problem, it has limitations. Our calculation relies on automated, single-selection, uncalibrated measurements from published literature, which may introduce biases and uncertainties. Additionally, we do not account for potential variations in xi\_ion or clumping factor C, which can significantly impact the results. Furthermore, our study does not incorporate new observational data, which could provide more accurate constraints on the parameters involved. These limitations highlight the need for further research and refined measurements to better understand the reionization process.
\section*{References}
\footnotesize
[Muoz2024] Muñoz et al. (2024). Reionization after JWST: a photon budget crisis?. bibcode 2024MNRAS.535L..37M \par [Park2022] Park et al. (2022). Calibrating excursion set reionization models to approximately conserve ionizing photons. bibcode 2022MNRAS.517..192P \par [Duncan2015] Duncan et al. (2015). Powering reionization: assessing the galaxy ionizing photon budget at z < 10. bibcode 2015MNRAS.451.2030D \par [Davies2021] Davies et al. (2021). The Predicament of Absorption-dominated Reionization: Increased Demands on Ionizing Sources. bibcode 2021ApJ...918L..35D \par [Madau2017] Madau (2017). Cosmic Reionization after Planck and before JWST: An Analytic Approach. bibcode 2017ApJ...851...50M

\end{document}
